\documentclass[
  spanish,
  letterpaper, oneside
]{article}

\def\documenttitle {Interaprendizaje en ambientes virtuales}
\def\documentsubtitle {Panorama de la asignatura}
\def\documentsubject {Actividad X - Interaprendizaje en ambientes virtuales}

\def\documentauthor {Martin Jonathan de la Cruz}
\def\coursename {Interaprendizaje en ambientes virtuales}
\def\coursecode {004}

\def\universityname {Universidad Abierta y a Distancia de Mexico}
\def\universityfaculty {Licenciatura en Derecho}
\def\universitydepartment {Interaprendizaje en ambientes virtuales}
\def\universitydepartmentimage {departamentos/UnADM}
\def\universitydepartmentimagecfg {height=1.57cm}
\def\universitylocation {Roma Norte, Ciudad de Mexico}

\def\authortable {
  \begin{tabular}{ll}
    \textbf{Alumno:}
    & \begin{tabular}[t]{l}
      Martin Jonathan de la Cruz \\
    \end{tabular} \\
    Matricula: & ES2611202040 \\
    Figura docente: & Janeth Santana Ramirez \\
    Semestre/Bloque: & 2 / 1 \\
    Tipo/Creditos: & Optativa / 6 \\
    & \\
    \multicolumn{2}{l}{Fecha de realizacion: \today} \\
    \multicolumn{2}{l}{\universitylocation}
  \end{tabular}
}

\input{template}
\usepackage{helvet}
\renewcommand{\familydefault}{\sfdefault}
\usepackage{array}
\usepackage{longtable}
\usepackage{pdflscape}
\usepackage{tikz}
\usetikzlibrary{arrows.meta,positioning,calc,matrix,fit,shapes.geometric,shadows.blur}
\setcitestyle{authoryear,open={(},close={)}}

\def\coverwatermarkenabled {true}
\def\coverwatermarkimage {img/departamentos/UnADM.pdf}
\def\coverwatermarkopacity {0.16}
\def\coverwatermarkwidth {0.72\paperwidth}
\def\coverwatermarkxshift {0cm}
\def\coverwatermarkyshift {-0.35cm}

\newcommand{\insertcoverwatermark}{%
  \ifthenelse{\equal{\coverwatermarkenabled}{true}}{%
    \AddToShipoutPictureBG*{%
      \begin{tikzpicture}[remember picture,overlay]
        \node[
          opacity=\coverwatermarkopacity,
          inner sep=0pt,
          xshift=\coverwatermarkxshift,
          yshift=\coverwatermarkyshift
        ] at (current page.center) {%
          \includegraphics[width=\coverwatermarkwidth]{\coverwatermarkimage}%
        };
      \end{tikzpicture}%
    }%
  }{}%
}

\newcommand{\pendiente}[1]{\textcolor{red}{[PENDIENTE: #1]}}

\begin{document}

\insertcoverwatermark
\templatePortrait
\templatePagecfg
\onehalfspacing

\begin{abstractd}
  Esta base editorial sintetiza la materia \textit{Interaprendizaje en ambientes virtuales}
  como una optativa de segundo semestre enfocada en aprender con otras personas,
  organizar trabajo juridico en linea y convertir la participacion digital en
  evidencia profesional verificable. Las notas internas funcionan como pivote de
  consigna y memoria de aula; las referencias visibles se sostienen en fuentes
  institucionales, normativas y academicas sobre aprendizaje colaborativo,
  competencias digitales, interculturalidad y acceso a derechos.
\end{abstractd}

\templateIndex
\templateFinalcfg

\section{Encuadre de la asignatura}

La materia propone una lectura no aislada del aprendizaje en linea. El eje no es
solamente cumplir actividades en una plataforma, sino construir conocimiento con
otras personas mediante dialogo, colaboracion, uso deliberado de TICCAD y
reflexion sobre diversidad cultural y linguistica. En clave juridica, esto se
traduce en practicar acuerdos de trabajo, trazabilidad de evidencia, escritura
colaborativa y coordinacion para una simulacion de despacho juridico digital.

El objetivo general es construir estrategias de interaprendizaje virtual que
permitan analizar como participamos, como nos comunicamos y como mejoramos la
autogestion, la coevaluacion y la produccion profesional en entornos digitales
orientados al Derecho. Esta orientacion es compatible con enfoques de aprendizaje
social y colaborativo que subrayan interaccion, responsabilidad individual e
interdependencia positiva \citep{vygotsky1978mind,johnsonJohnson1999learning}.

\section{Proposito formativo}

La secuencia de la asignatura combina reflexion, practica y aplicacion. Primero
se reconoce la diversidad del grupo y el sentido del interaprendizaje; despues
se define un ecosistema TICCAD minimo y funcional; finalmente, el curso culmina
en un proyecto colaborativo con logica profesional: un despacho juridico digital
que organiza roles, acuerdos, herramientas y una propuesta vinculada con un caso.

La materia remarca cuatro expectativas permanentes: autonomia, participacion
proactiva, respeto en la interaccion y aprovechamiento reflexivo de las
herramientas digitales. Esas expectativas deben notarse en cada reporte,
presentacion y producto visual de la materia, con apoyo en competencias digitales
e interculturales verificables \citep{unescoDigitalLiteracy2018,unescoInterculturalCompetences2013}.

\section{Ruta de evaluacion}

\begin{longtable}{p{1.2cm}p{9.0cm}p{2.0cm}}
\textbf{Semana} & \textbf{Actividad} & \textbf{Valor} \\
\hline
1 & Diagnostica & 0/100 \\
2 & Foro: diversidad cultural y linguistica en el Derecho & 10/100 \\
3 & Interaprendizaje: comprender, construir y distinguir & 10/100 \\
4 & Mi compromiso personal como estudiante en entornos virtuales & 0/100 \\
5 & Construccion de un ecosistema TICCAD para el trabajo juridico virtual & 15/100 \\
6 & Documento colaborativo: organizacion del despacho juridico digital & 20/100 \\
7 & Brief profesional del despacho juridico digital & 0/100 \\
8 & Propuesta profesional del despacho juridico digital & 30/100 \\
9 & Bitacora de interaprendizaje & 15/100 \\
10 & Autoevaluacion & 0/100 \\
\hline
 & \textbf{Total} & \textbf{100/100} \\
\end{longtable}

La ponderacion muestra que el curso valora mas la capacidad de coordinarse,
documentar decisiones y presentar una propuesta juridica coherente que la mera
presencia formal en el aula. Las actividades sin puntaje directo siguen siendo
relevantes porque preparan la colaboracion y la reflexion final.

\section{Pauta de realizacion}

Toda entrega debe transformar la consigna en un problema de organizacion,
participacion o produccion juridica virtual claramente delimitado. La redaccion
debe incluir contexto, conceptos, criterio de seleccion de herramientas,
evidencia de colaboracion, analisis propio y una consecuencia transferible a la
practica profesional.

\begin{itemize}
  \item La diversidad cultural y linguistica se trata como condicion real del trabajo en equipo, no como nota decorativa.
  \item Las TICCAD se eligen por funcion, accesibilidad y trazabilidad.
  \item El proyecto colaborativo debe operar con acuerdos verificables de roles, tiempos y control de versiones.
  \item La autoevaluacion y la coevaluacion solo son validas si se sustentan en hechos observables.
\end{itemize}

\section{Estructura sugerida}

\begin{enumerate}
  \item Encuadre: que reto de interaprendizaje virtual aborda la actividad y por que importa en la formacion juridica.
  \item Desarrollo: conceptos de colaboracion, diversidad, TICCAD o gestion de trabajo aplicados a la consigna.
  \item Evidencia: cuadro, protocolo, matriz, bitacora, mapa o pieza profesional que haga visible el proceso.
  \item Valoracion: que funciono, que friccion aparecio y que decision mejora la coordinacion futura.
  \item Cierre: aprendizaje transferible a un entorno juridico digital real.
\end{enumerate}

\section{Checklist editorial}

\begin{itemize}
  \item La actividad responde a la planeacion sin copiar de forma mecanica la consigna.
  \item Las fuentes citadas aparecen en la bibliografia local y se usan con un proposito claro.
  \item El producto visual o colaborativo deja rastro de acuerdos, roles o decisiones de trabajo.
  \item La identidad UnADM aparece en portada, tono academico e integridad de la entrega.
  \item La conclusion no se limita a resumir: fija una mejora concreta para el trabajo juridico virtual.
\end{itemize}

\section{Criterios de entrega y evaluacion}

La materia exige distinguir entre participar y construir. Participar es aparecer;
interaprender es modificar la comprension propia a partir del trabajo con otras
personas. Por eso, cada actividad debe mostrar alguna forma de interdependencia:
acuerdos previos, evidencia de revision cruzada, uso justificado de herramientas,
registro de decisiones o reflexion sobre como incidio la diversidad del grupo en
el resultado final.

\section{Conclusion editorial}

Interaprendizaje en ambientes virtuales funciona como una antesala a la practica
profesional juridica en linea. Su valor no radica en acumular plataformas o en
simular colaboracion, sino en aprender a coordinar criterios, documentar
evidencia, sostener dialogos respetuosos y convertir la organizacion digital en
un producto juridico comprensible. Esa es la regla de lectura que debe orientar
las futuras actividades de esta carpeta.

\clearpage
\nocite{unadmSitioWeb,unadmMallaDerecho2024,vygotsky1978mind,johnsonJohnson1999learning,unescoDigitalLiteracy2018,unescoInterculturalCompetences2013,unescoGenerativeAI2023,cpeum2026,leyDerechosLinguisticos2024}
\bibliography{interaprendizaje-en-ambientes-virtuales}

\end{document}